\documentclass[11pt]{article}
\usepackage[a4paper,margin=26mm]{geometry}
\usepackage{amsmath,amssymb,booktabs,longtable,array}
\usepackage[unicode,hidelinks]{hyperref}
\usepackage{fontspec}
\setlength{\emergencystretch}{2em}
\setmainfont{lmroman10-regular.otf}[BoldFont=lmroman10-bold.otf,ItalicFont=lmroman10-italic.otf,BoldItalicFont=lmroman10-bolditalic.otf]
\newcommand{\tightlist}{\setlength{\itemsep}{0pt}\setlength{\parskip}{0pt}}
\title{Marked coordinates and exact orbit enumeration}
\author{human\_compression research project}
\date{11 October 2026 --- author manuscript}
\begin{document}\maketitle
\section{Coordinates on a marked
group}\label{coordinates-on-a-marked-group}

A marking is an ordered tuple, not an unordered generating set. The
distinction matters whenever the positions specify face ranks. Let \(G\)
be a group and let \(\boldsymbol s=(s_1,\ldots,s_m)\) be a tuple in
\(G\). Define \[
a_0=1,\qquad a_i=s_1\cdots s_i\quad(1\leq i\leq m).
\tag{1}
\] The reverse substitution is \(s_i=a_{i-1}^{-1}a_i\). For a set
\(J\subseteq\{0,\ldots,m\}\) put \[
 D_J(\boldsymbol a)=\langle a_i^{-1}a_j:i,j\in J\rangle,
\quad
 T_J(\boldsymbol s)=\langle s_{i+1}\cdots s_j:i<j,\ i,j\in J\rangle.
\tag{2}
\] An empty generating set generates the trivial subgroup.

\textbf{Theorem MC.1 (coordinate transport).} The substitutions in (1)
are mutually inverse, preserve the generated group, commute with every
group homomorphism, and satisfy \[
a_i^{-1}a_j=s_{i+1}\cdots s_j\quad(i<j),\qquad
D_J(\boldsymbol a)=T_J(\boldsymbol s).
\tag{3}
\] Consequently every assertion involving these particular subgroup
diagrams, their intersections, their generated joins and their element
orders can be transferred by substitution, with no new group-theoretic
hypothesis.

\emph{Proof.} Adjacent factors cancel in
\((a_i^{-1}a_{i+1})\cdots(a_{j-1}^{-1}a_j)\), leaving \(a_i^{-1}a_j\).
Conversely, cancellation in \((s_1\cdots s_i)^{-1}(s_1\cdots s_j)\)
gives the first equality in (3). Each original generator is a word in
the new ones and conversely, so generation is preserved. A homomorphism
commutes with products and inverses. In (2), a difference with reversed
endpoints is the inverse of the difference with increasing endpoints,
and equal endpoints contribute the identity. Thus the two displayed
generating sets generate exactly the same subgroup. All the asserted
transports now follow by replacing equal subgroups or equal elements.
\(\square\)

This elementary substitution is the marked-coordinate principle already
recorded as M2 in the project proof bank. Its application to chiral
polytopes is the substitution of Buekenhout--Leemans--Tranchida, Theorem
4.3; neither that substitution nor its use of \(C^+\)-groups is claimed
as new here. The formal certificate accompanying this paper verifies (3)
for arbitrary interval lengths and the equality of the actual subgroup
closures, including reversed endpoint pairs.

\textbf{Lemma MC.2 (a star generates the difference subgroup).} For
nonempty \(J\) and \(b\in J\), \[
D_J(\boldsymbol a)=\langle a_b^{-1}a_j:j\in J\setminus\{b\}\rangle.
\tag{4}
\]

\emph{Proof.} The right side is contained in the left. Every generator
of the left has the expression
\(a_i^{-1}a_j=(a_b^{-1}a_i)^{-1}(a_b^{-1}a_j)\), proving the reverse
inclusion. \(\square\)

Formula (4) reduces a subgroup computation with quadratically many named
differences to one with at most \(|J|-1\) generators. It changes neither
the subgroup nor the intersection tests.

\section{Exact equivalence of generating
markings}\label{exact-equivalence-of-generating-markings}

We recall the shared orbit interface used in the enumeration topic. If a
group \(B\) acts on sets \(Z,E\), and \(p:Z\to E\) is equivariant, then
\[
 B\backslash Z\ \simeq\ \coprod_{e\in B\backslash E}B_e\backslash p^{-1}(e),
\tag{5}
\] where one base representative is chosen in each orbit. This is
U.ExtensionOrbits. The proof and its previous geometric applications are
retained from the shared umbrella: first move the base into the selected
orbit representative; any remaining equivalence fixes that base. Formula
(5) applies to a tuple and its initial subtuple. In particular, after
fixing one entry one must use its stabilizer, rather than continuing to
quotient the remaining entries independently by the whole group.

For a finite group \(G\), let \(\mathcal M_m(G)\) be the set of ordered
generating \(m\)-tuples. The automorphism group acts freely on this set:
an automorphism fixing all the generators fixes every word. Thus every
orbit has size \(|\operatorname{Aut}(G)|\). One need not construct the
whole automorphism group to test this equivalence.

\textbf{Theorem MC.3 (Cayley transport).} Let \(\boldsymbol a\) and
\(\boldsymbol b\) be generating \(m\)-tuples of finite groups \(G,H\) of
equal order. The assignment \(a_i\mapsto b_i\) extends to a group
isomorphism if and only if the following finite propagation succeeds.
Start with \(f(1)=1\), and whenever \(f(x)\) is assigned, require \[
f(xa_i)=f(x)b_i\quad(1\leq i\leq m).
\tag{6}
\] Reject if two requirements disagree. Accept if the resulting
assignment is a bijection. The isomorphism, when it exists, is unique.

\emph{Proof.} Every group isomorphism with the required generator images
satisfies (6), so disagreement or failure of bijectivity is a valid
rejection. A finite group generated as a group by the \(a_i\) is also
generated as a monoid: \(a_i^{-1}\) is a positive power of \(a_i\).
Hence propagation visits all of \(G\). If it succeeds, induction on a
positive word \(w\) gives \(f(xw(\boldsymbol a))=f(x)w(\boldsymbol b)\).
Taking \(x=1\) shows that the right side depends only on the element
represented by \(w\). Taking a word for an arbitrary \(y\) then gives
\(f(xy)=f(x)f(y)\). Acceptance therefore gives a bijective homomorphism.
Uniqueness follows from generation. \(\square\)

A canonical implementation starts at the identity, visits new elements
in breadth-first order using the ordered generators, and records the
target numbers of all colored Cayley edges. Two such finite arrays agree
exactly when the markings are isomorphic. Indeed equality of arrays
supplies the bijection in MC.3; conversely an isomorphism carries the
deterministic traversal to the same traversal. This test does not
quotient the tuple by a permutation of its positions.

The same method applies to connected graphs with one edge of each color
at every vertex: a color-preserving graph map is forced by the image of
one vertex. Trying every image, rejecting inconsistencies and testing
bijectivity enumerates the full color-preserving automorphism group.
This is a direct finite certificate, not a heuristic based on a chosen
subgroup of graph automorphisms.

\section{Hereditary searches and
termination}\label{hereditary-searches-and-termination}

Suppose an admissibility predicate is defined on normalized tuples
\((1,a_1,\ldots,a_m)\) and has the following properties:

\begin{enumerate}
\def\labelenumi{\arabic{enumi}.}
\tightlist
\item
  every initial subtuple is admissible in the subgroup it generates;
\item
  admissibility is preserved by automorphisms of the ambient group;
\item
  every appended entry lies in a prescribed finite candidate set;
\item
  every admissible extension enlarges the subgroup generated by the
  prefix.
\end{enumerate}

\textbf{Theorem MC.4 (exhaustive prefix search).} Start at
representatives of the allowed first-entry conjugacy classes.
Recursively append every candidate that passes the admissibility tests.
Record a tuple when it generates the input group, and stop that branch.
Deduplication by MC.3 gives exactly the automorphism orbits of all
admissible generating tuples.

\emph{Proof.} Every retained tuple is admissible by the tests, and is
recorded only when it generates the required group. Conversely,
conjugate any admissible generating tuple to put its first entry into
the chosen conjugacy representative. Every prefix passes the tests by
hypothesis 1; induction on its length shows that the recursion reaches
the entire tuple. Hypothesis 4 shows that no proper prefix already
generates the ambient group, so stopping a branch early loses no output.
The search terminates because each subgroup index along a strict finite
subgroup chain is at least two. Finally MC.3 removes exactly, and all,
the automorphism equivalences among the recorded tuples. \(\square\)

The conjugacy reduction need not remove every automorphism equivalence
at the first entry. Inner automorphisms suffice for completeness, and
final automorphism deduplication removes any remaining duplicates.
Alternatively, (5) permits successive reduction by the actual stabilizer
of each prefix. Such reduction is optional; its implementation must
preserve the prefix rather than treating entries independently.

A useful sufficient condition for strict growth is the
difference-subgroup intersection property \[
 D_J\cap D_K=D_{J\cap K}.
\tag{7}
\] For a prefix \((1,a_1,\ldots,a_{k-1})\) and a nonidentity new entry
\(a_k\), take \(J=\{0,\ldots,k-1\}\) and \(K=\{0,k\}\). Then
\(\langle a_1,\ldots,a_{k-1}\rangle\cap\langle a_k\rangle=1\). In
particular \(a_k\) is not in the old subgroup. If \(|a_1|\geq3\) and the
final tuple has \(n\) coordinates including \(a_0\), then \[
 |G|\geq 3\,2^{n-2}.
\tag{8}
\] The implementation may use this bound or simply stop whenever
generation occurs. Its proof of termination is independent of a
numerical cutoff selected by the programmer.

\section{Reflection tests}\label{reflection-tests}

Let \((1,a_1,\ldots,a_{n-1})\) generate \(G\), assume \(a_i^2=1\) for
\(i\geq2\), and write \(s_i=a_{i-1}^{-1}a_i\).

\textbf{Proposition MC.5 (reflection transport).} An automorphism \(f\)
satisfies \[
f(a_i)=a_i^{-1}\quad(1\leq i<n)
\tag{9}
\] if and only if \[
f(s_1)=s_1^{-1},\qquad f(s_2)=s_1^2s_2,
\qquad f(s_i)=s_i\quad(i\geq3).
\tag{10}
\] Any such automorphism has square equal to the identity. If
\(|a_1|>2\), it is not the identity.

\emph{Proof.} Condition (9) fixes \(a_i\) for \(i\geq2\). Therefore
\(f(s_2)=f(a_1^{-1}a_2)=a_1a_2=a_1^2s_2\), and both factors in
\(s_i=a_{i-1}^{-1}a_i\) are fixed for \(i\geq3\). This gives (10).
Conversely (10) sends the first prefix to its inverse and fixes the
second prefix \(s_1s_2\). All subsequent prefixes are then fixed. These
are exactly (9), since they are involutions. Finally \(f^2\) fixes every
generating entry, and hence is the identity. Its action on \(a_1\) is
nontrivial when \(|a_1|>2\). \(\square\)

This is the content of BLT Lemma 4.4 in the correct automorphism
language, expressed through MC.1. It concerns all group automorphisms,
not only conjugations by elements of \(G\). MC.3 decides (9) without
requiring the automorphism group to be enumerated in advance.

\section{Attribution, formal coverage and
reuse}\label{attribution-formal-coverage-and-reuse}

The umbrella collects elementary coordinate and enumeration interfaces.
MC.1 is the existing marked-coordinate principle; the specific prefix
substitution and reflection transport used here belong to BLT. The
orbit-over-a-base result (5) is U.ExtensionOrbits, already used in the
accepted enumeration module. MC.2, MC.3 and MC.4 supply explicit
complete proofs of the algebraic and algorithmic details needed by
applications; no priority is asserted for these elementary statements.

\texttt{PrefixCoordinates.lean} proves seven declarations:
arbitrary-length telescoping, recovery in both coordinate directions,
equality of subgroup closures, transport of the full intersection
assertion, the reflected second generator, and the square of an
inversion automorphism on an arbitrary generating set. The subgroup
theorem is not merely equality of two generated groups at the top: it
identifies every subgroup indexed by a set of endpoints. The formal file
does not prove the geometric reconstruction theorem, finite search
completeness, the involution-count theorem for almost simple groups, or
the correctness of the Python interpreter and GAP.

The associated application must pay for its geometric recognition
theorem and for its specific admissibility conditions. This umbrella is
charged in full to any topic using it under the project's requested
topic-local measurement convention. Shared acquisition inside one topic
is recorded once; no numerical compression ratio is claimed here.

\end{document}
